\section{Contact sensitivity and Newton conditioning}
\label{sec:conditioning-counterexamples}

The approximation obstructions concern derivatives of selected contact
maps across labels.  A Newton condition number concerns a linear system
at one interior point.  The following two separations identify why an
additional sensitivity theorem is needed to connect them.

\subsection{A sparse change of cone coordinates}

Suppose $F(Gz)=F(z)+c_G$ for a linear cone automorphism $G$.  In an
equality formulation $Ez=b$, put $z'=Gz$, $E'=EG^{-1}$ and
$H=\nabla^2F(z)$.  Differentiation gives
\begin{equation}\label{eq:conditioning-schur}
 H'=G^{-T}HG^{-1},\qquad
 E'(H')^{-1}(E')^T=EH^{-1}E^T.
\end{equation}
The raw augmented system instead changes by congruence:
\[
 \begin{pmatrix}H'&(E')^T\\E'&0\end{pmatrix}
 =\begin{pmatrix}G^{-T}&0\\0&I\end{pmatrix}
  \begin{pmatrix}H&E^T\\E&0\end{pmatrix}
  \begin{pmatrix}G^{-1}&0\\0&I\end{pmatrix}.
\]
Thus the Schur matrix is literally unchanged, whereas the Euclidean
spectrum of the augmented matrix need not be.

For $Q_m$, $m\geq3$, let $a=e_0+e_1$, $b=e_0-e_1$, and let $G_\lambda$
be the identity on transverse coordinates, with first block
\[
 \frac12\begin{pmatrix}
 \lambda+\lambda^{-1}&\lambda-\lambda^{-1}\\
 \lambda-\lambda^{-1}&\lambda+\lambda^{-1}
 \end{pmatrix},\qquad \lambda>0.
\]
It preserves the Lorentz quadratic form and its future sheet.  Moreover,
$G_\lambda a=\lambda a$ and $G_\lambda^{-T}b=\lambda b$.
Consequently both primal and dual Euclidean norms can approach zero while
all their pairings remain unchanged.  This is a genuine contact example:
\[
 A(u)=(\sqrt{1+\norm u^2},1,u),\qquad
 B(v)=(\sqrt{1+\norm v^2},-1,-v)
\]
satisfy $\ip{A(u)}{B(v)}\geq0$, with equality at $u=v$.
At zero their transverse derivatives are isometries and are unchanged by
the boost; the base norms scale by $\lambda$.
The standard Lorentz barrier is invariant, so
\eqref{eq:conditioning-schur} applies.  At $e_0$, $H=2I$, but at
$G_\lambda e_0$ it is $2G_\lambda^{-2}$, with condition number
$\lambda^{-4}$ for $0<\lambda<1$.
This calculation prevents a lower bound involving unnormalized Euclidean
factor norms from being read as a lower bound on the intrinsic Schur
condition number.  The boost only mixes two columns of $E$ but can change
their magnitudes, data-access normalization, and the cost of a solver
using the raw augmented system.

\subsection{Bounded pointwise geometry and unbounded contact derivatives}

Transporting the normalization covector avoids the preceding coordinate
issue, but leaves a second obstruction.  For maps $A,B$ and a strictly
positive covector $\ell$, the transformations
\[
 (A,B,\ell)\mapsto(GA,G^{-T}B,G^{-T}\ell)
\]
preserve $\ip AB$, $-dA^*dB$, $A^*dB$, and $d\log\ell(A)$.
Their cross-label derivatives can nevertheless grow while all pointwise
interior geometry stays bounded.

Use the Lorentz Jordan product
$(t,z)\circ(s,w)=(ts+z^Tw,tw+sz)$ and fix
$c_\pm=(1,\pm p_0)/2$ for $p_0\in S^1$.
On $S^n$, $n\geq1$, put $a_k(x)=\exp(\sin(kx_1))$ and
\begin{equation}\label{eq:conditioning-radial}
 A_k=a_kc_+,\quad B_k=a_k^{-1}c_-,\qquad
 X_k=a_k(c_++\tau c_-),\quad
 S_k=a_k^{-1}(c_-+\tau c_+),
\end{equation}
where $0<\tau<1$ is fixed.  The boundary maps are complementary,
$e^{-1}\leq a_k\leq e$, and
$\norm{d\log\ell(A_k)}_\infty=k$ for any positive $\ell$.
Yet $X_k\circ S_k=\tau e$ exactly.  Since
$\nabla^2F(az)=a^{-2}\nabla^2F(z)$, the primal and dual Hessian norms
and inverse norms are uniformly bounded in $k$.  In the common Jordan
frame the Nesterov--Todd scaling point has eigenvalues
$a_k/\sqrt\tau$ and $a_k\sqrt\tau$; its spectral ratio is $1/\tau$.
Here centrality uses the displayed Jordan-product convention; a constant
trace-inner-product normalization has no effect on the conclusion.

For a different example fix $0<\theta<\pi/2$ and put
$p_k(x)=(\cos(kx_1),\sin(kx_1))$, $q_k=R_\theta p_k$, and
\begin{equation}\label{eq:conditioning-rotating}
 A_k=(1,p_k),\quad B_k=(1,-q_k),\qquad
 X_k=(1+\tau,p_k),\quad S_k=(1+\tau,-q_k).
\end{equation}
The diagonal defect is the constant $1-\cos\theta$, but
\[
 d\ip{A_k}{B_k}=0,\qquad
 A_k^*dB_k=\sin\theta\,d(kx_1),\qquad
 \norm{A_k^*dB_k}_\infty=k\sin\theta.
\]
All interior pairs are common spatial rotations of one fixed pair, so
Hessian spectra and scaling-point spectra are independent of $k$.
Their relative Jordan central residual is
\[
 \frac{2(1+\tau)\sin(\theta/2)}
 {2\tau+\tau^2+1-\cos\theta}.
\]
For any prescribed positive central-neighborhood radius, sufficiently
small fixed $\theta>0$ puts the entire family inside that neighborhood.
Alternatively, $\theta=k^{-1/2}$ gives defect $\Theta(k^{-1})$ and
central residual $O(k^{-1/2})$, while the one-sided derivative grows as
$\Theta(\sqrt{k})$.  The mixed channel is
$\cos\theta\,d(kx_1)\otimes d(kx_1)$; it vanishes at the critical
points of $x_1$, consistently with the phase obstruction.

These explicit computations prove that pointwise centrality, Hessian
conditioning, scaling-point conditioning, and bounded factor scales do
not control either radial or one-sided contact sensitivity.  The
pointwise interpretation of self-scaled geometry is classical
\citep{NT1997}; the examples distinguish it from the different,
cross-objective quantities used here.
One may even impose the identity equality $z=X_k$ at each label.  The
assembled saddle matrix $\left(\begin{smallmatrix}\mu H&I\\I&0
\end{smallmatrix}\right)$ and its inverse are uniformly bounded whenever
$\mu$ stays in a fixed positive compact interval, since its inverse is
$\left(\begin{smallmatrix}0&I\\I&-\mu H\end{smallmatrix}\right)$.
These artificial instances have changing right-hand sides.  They are
not a single approximate lift of a fixed body.

\subsection{What an additional sensitivity hypothesis would give}

In the robust submersion inequality~\eqref{eq:submersion-threshold}, suppose the finite-radius estimate
reduces to
\[
 \delta+K_A\sqrt{2H_B\epsilon}\geq\mu,
 \qquad \sqrt{2\epsilon/H_B}\leq r_0,
\]
where $\delta<\mu$, $\epsilon>0$ and $H_B>0$.
Then
$H_B\geq(\mu-\delta)^2/(2K_A^2\epsilon)$.
The finite-radius expression $\inf_{0<s\leq r_0}
(\epsilon/s+H_Bs/2)$ remains the valid form when $H_B=0$ or the
optimizing radius does not fit; one cannot divide by $H_B$ in those cases.
If a separately established, representation-specific sensitivity theorem
provides $K_A\leq c_K\chi^a$ and $H_B\leq c_H\chi^b$, with positive
constants and exponents, substitution yields
\[
 \chi\geq
 \left(\frac{\mu-\delta}{c_K\sqrt{2c_H\epsilon}}\right)^{1/(a+b/2)}.
\]
Likewise, if $C>0$, $K_0>0$, and $q>0$, then
$H_B\leq C2^{q\beta}$ and $K_A\leq K_0$ imply
\[
 \beta\geq\frac1q\log_2
 \frac{(\mu-\delta)^2}{2CK_0^2\epsilon}.
\]
These are conditional algebraic consequences.  Neither an ordinary
Newton condition number nor a coefficient bit count is proved to satisfy
these extra hypotheses.  In particular, differentiating a KKT solution
map requires control of both its inverse Jacobian and the variation of
the input data with the contact label.
